\section{Dependency Chain}
\label{sec:chain}

The driver links two libraries.
\texttt{algebra::matrices} stores the entries and implements matrix multiplication and row swaps.
\texttt{algebra::finite\_fields} supplies the entry type.
The driver fixes that type as \(\mathtt{FiniteField}\langle 11,3\rangle\), so every entry is an element of a degree-3 extension of the prime field \(\mathbb{F}_{11}\) \cite{russell2026fields,russell2026matrices}.
Polynomial arithmetic used by the field lives in a third library, pulled in by the finite-field package.
The matrices never reduce polynomials themselves.
They call \texttt{add}, \texttt{multiply}, \texttt{negate}, and \texttt{reciprocal} on the element type.

That split is why the finite-field element is the class to cache.
A change inside the element type is visible to multiplication and to elimination without editing the matrix templates or the driver.
The irreducible modulus stays the one chosen by the field's deterministic search.
On the profiled build the modulus was \(x^3 + x + 4\).

The sequential driver still computes \(C = AB\) with the library matrix product, then runs Gaussian elimination with partial pivoting for \(\det(A)\), \(\det(B)\), and \(\det(C)\).
The elimination is local to the driver.
It is the same algorithm for the original field and the cached field, so a timing difference between those two builds is a difference in element arithmetic.
